\subsection{SUBGROUPS AND QUOTIENT GROUPS OF A QUOTIENT GROUP}

Let G be a topological group, let H be a normal subgroup of G, and let \(\varphi: G \to G/H\) be the canonical mapping. We know that if \(A'\) is a subgroup of G/H, then \(\varphi^{-1}(A')\) is a subgroup of G which contains H. Conversely, if A is a subgroup of G, then \(\varphi(A)\) is a subgroup of G/H; moreover, there is a canonical bijection of the quotient group \(A/(A \cap H)\) onto the subgroup \(\varphi(A)\) of G/H, and a canonical bijection of \(\varphi(A)\) onto the quotient group AH/H, and both these bijections are isomorphisms for the group structures.

PROPOSITION 20. Let A be a subgroup of a topological group G, let H be a normal subgroup of G, and let \(\varphi\) denote the canonical mapping of G onto G/H. Then the canonical bijection of \(\varphi(A)\) onto AH/H is an isomorphism of topological groups.

This follows from the preceding remarks, and from Proposition 10 of no. 4.

\pandocbounded{\includegraphics[keepaspectratio]{images/part002_0097f9c2.jpg}}

The canonical bijection of \(\mathbf{A}/(\mathbf{A} \cap \mathbf{H})\) onto \(\varphi(\mathbf{A})\) is a continuous homomorphism, since it arises from the restriction of \(\varphi\) to \(\mathbf{A}\) by passing to the quotients; but in general the topological groups \(\mathbf{A}/(\mathbf{A} \cap \mathbf{H})\) and \(\mathbf{AH}/\mathbf{H}\) are not isomorphic (cf.~§ 4, no. 1, Proposition 1, Corollary 3).

* For example, take G to be the additive group R of real numbers, take H to be the group Z of integers, and take A to be the group θZ of integer multiples of an irrational number θ. Then A∩H = \{0\}, so that A/(A∩H) is a discrete group, isomorphic to Z; on the other hand, A + H is dense in R (as we shall see in Chapter V, § 1, no. 1, Proposition 1), hence (A + H)/H, which is locally isomorphic to A + H (no. 6, Proposition 19), is not a discrete group and consequently is not isomorphic to A/(A∩H). *

Nevertheless, we have the following proposition:

PROPOSITION 21. Let G be a topological group, \(G_{0}\) a dense subgroup of G, \(H_{0}\) a normal subgroup of \(G_{0}\) , H the closure of \(H_{0}\) in G and \(\varphi\) the canonical mapping \(G \rightarrow G/H\) . Then the canonical bijection \(G_{0}/H_{0} \rightarrow \varphi(G_{0})\) is an isomorphism of the topological group \(G_{0}/H_{0}\) onto a dense subgroup of G/H.

Since \(\mathbf{H}_0 = \mathbf{H} \cap \mathbf{G}_0\) , it is enough to show that if \(\mathbf{U}_0\) is any open subset of \(\mathbf{G}_0\) which is saturated with respect to the relation \(x^{-1}y \in \mathbf{H}_0\) , then \(\mathbf{U}_0\) is the intersection of \(\mathbf{G}_0\) and an open subset of \(\mathbf{G}\) which is saturated with respect to the relation \(x^{-1}y \in \mathbf{H}\) (Chapter I, § 3, no. 6, Proposition 10). Let \(\mathbf{U}\) be an open subset of \(\mathbf{G}\) such that \(\mathbf{U}_0 = \mathbf{U} \cap \mathbf{G}_0\) . Since \(\mathbf{U}_0 = \mathbf{U}_0\mathbf{H}_0\) , it is easily seen that \(\mathbf{U}_0 = \mathrm{UH}_0 \cap \mathbf{G}_0\) ; but \(\mathrm{UH}_0\) is open in \(\mathbf{G}\) , so we may suppose that \(\mathbf{U} = \mathrm{UH}_0\) . The set \(\mathrm{UH}\) is open in \(\mathbf{G}\) and is saturated with respect to the relation \(x^{-1}y \in \mathbf{H}\) ; hence the proposition will be proved if we show that \(\mathrm{UH} \cap \mathbf{G}_0 = \mathbf{U}_0\) . Now, if \(u \in \mathbf{U}\) and \(h \in \mathbf{H}\) are such that \(uh \in \mathbf{G}_0\) , then there is a symmetric neighbourhood \(\mathbf{V}\) of \(e\) in \(\mathbf{G}\) such that \(uv \subset U\) ; since \(Vh\) is a neighbourhood of \(h\) in \(\mathbf{G}\) , there exists \(z \in V\) such that \(zh \in \mathbf{H}_0\) . But then we have \(uz^{-1} \in U\) , and \(uh = (uz^{-1})(zh)\) ; therefore \(\mathrm{UH} \cap \mathbf{G}_0 \subset \mathrm{UH}_0\) . But \(\mathrm{UH}_0 = U\) , hence \(\mathrm{UH} \cap \mathbf{G}_0 \subset U \cap \mathbf{G}_0 = U_0\) . This completes the proof.

Let G be a topological group operating continuously on a topological space X, and let K be a normal subgroup of G which is contained in the stabilizer of each point of X. For each \(x \in X\) , the relation \(s \equiv t \pmod{K}\) therefore implies s.x = t.x, and passing to the quotient we define a mapping \(\dot{s} \to \dot{s}.x\) of G/K into X. It is immediately verified that with respect to the external law \((\dot{s}, x) \to \dot{s}.x\) , X has G/K as a group of operators. Moreover, G/K operates continuously on X with respect to this law; for the relation of equality on X and the relation \(s \equiv t \pmod{K}\) on G are both open equivalence relations, and therefore the result follows from the continuity of the mapping

\[
(s, x) \rightarrow \dot {s}. x = s. x
\]

of \(\mathbf{G} \times \mathbf{X}\) into \(\mathbf{X}\) (Chapter I, § 5, no. 3, Corollary to Proposition 8, and § 3, no. 4, Proposition 6).

Now let G be a topological group operating continuously on a topological space X, and let H be any normal subgroup of G; then H operates continuously on X. Let S be the equivalence relation defined by H on X; then S is open (no. 4, Lemma 2). The relation S is compatible with the group G (no. 4); for if \(y \equiv x \pmod{S}\) then there exists \(t \in H\) such that y = t.x; hence for all \(s \in G\) we have \(s.y = (sts^{-1}) \cdot s.x\) , and \(sts^{-1} \in H\) since H is normal in G; thus \(s.y \equiv s.x \pmod{S}\) . If \(\psi\) is the canonical mapping of X onto X/S, the group G therefore operates continuously on X/S with respect to the external law

\[
(s, \psi (x)) \rightarrow \psi (s. x)
\]

(no. 4, Proposition 11). Moreover, the group H is contained in the stabilizer of each of the points of X/S; as we have seen above, G/H operates continuously on X/S = X/H with respect to the external law \((s, \psi(x)) \to \psi(s.x)\) . If R denotes the equivalence relation on X defined by G, then the relation S implies R, and the equivalence relation R/S on X/S is that defined by the group G/H. Hence (Chapter I, § 3, no. 4, Proposition 7):

PROPOSITION 22. Let G be a topological group operating continuously on a topological space X, and let H be a normal subgroup of G. Then the canonical bijection of X/G onto (X/H)/(G/H) is a homeomorphism.

COROLLARY. Let G be a topological group, H a normal subgroup of G, and K a normal subgroup of G which contains H. Then the canonical bijection of G/K onto (G/H)/(K/H) is an isomorphism of topological groups.

We know already that this bijection is an isomorphism of groups, and Proposition 22 (applied to the group K operating on the right on G) shows that it is a homeomorphism.

